\subsection{Preserving the order of eligibility decisions}
\label{sec:ordered-event-eligibility}

A customer can be unknown before a collection is authorised. Suppose the full
intent supplies an unknown customer, false approval and a claim key absent
from the book. The retained lifecycle first reports unresolved input. Once the
customer is supplied, the same false approval produces an unsupported event;
only after approval can the missing claim produce an invalid event. Combining
all three predicates into one conjunction changes this behaviour: an invalid
consequence can outrank the earlier unresolved or unsupported result. The
individual predicates remain meaningful, but their conjunction computes a
different decision.

\code{OrderedEventEligibility/1} makes this order explicit. A template declares
its company context, operation, complete input types and ordered stages. Each
stage supplies one existing guarded requirements manifest, exact bindings for
its declared inputs, and the reason to report if it does not pass. The evaluator
validates the typed book, event identity, chronology, full input population,
types, values and provenance before interpreting the stages. Any unknown full
input then returns unresolved, preserving the input names. Approval is the
first semantic stage in each retained template. A stage with no active
requirement is refused rather than treated as permission to continue.

For a sale or collection, the next stage checks that assessment has not sealed
the event date. A sale then checks the due date and absence of the claim key.
A collection checks presence of the key before comparing its recorded customer
with the supplied customer. An assessment instead checks its starting cursor,
nonempty interval and last included date, then its period boundary. Execution
stops at the first stage that does not pass, retaining its status, reason and
stage identifier. The existing guarded evaluator still determines each typed
predicate; the ordered layer introduces no competing Boolean language.

The bindings supply facts rather than conclusions. They can read a full-intent
input, event identifier or date, book date or period end, a declared memory
leaf, map membership, or a selected row's scalar field. Membership is computed
against the actual typed map using the supplied text key. A caller cannot
replace that computation with an asserted Boolean. A row-field binding names
an earlier stage whose unconditional predicate establishes membership in the
same map under the same key. Compilation checks that relationship and the
declared field type without requiring a row to exist. It also checks every
later stage and guarded branch before execution, including when the map is
empty or an earlier approval will fail.

Calendar adjacency is a narrow additional binding. After strict date
validation, it compares the two date ordinals and returns true when the right
date is one calendar day after the left. An assessment therefore relates its
event date to its explicit end-exclusive date without constructing a successor
that could overflow at the maximum representable date. Leap days and year
changes follow the same relation. This is a civil-calendar condition, not a
business-day adjustment or a fiscal-calendar policy. The unchanged chronology
and interval checks remain separate conditions.

The call binds the template, book, event, schema registry and memory definition.
It also binds the order of full-intent inputs. Canonical JSON identity ignores
map insertion order, whereas the retained unresolved result lists unknown
inputs in that order. Reversing two unknown inputs therefore leaves the
canonical event hash unchanged but changes the order digest and the returned
list. Reusing the earlier call is refused. This is the same distinction between
value identity and explicitly meaningful order that arose for assessment writes.

Six templates across the two retained charts admit all 38 previously checked
events. Combined with the parameter and ordered-row constructors, they reproduce
the complete saved requests and transitions while native lifecycle eligibility
and construction functions are disabled. The original 18 refusals and 16
additional precedence cases preserve their statuses and reasons. Seventy
controls include selected-row binding, skipped-stage compilation, changed rule
truth, stale identities and calendar boundaries. A changed customer predicate
can admit a wrong customer; that consequential edit demonstrates author control
and deliberately differs from the selected contract.

These comparisons establish a finite declarative alternative for eligibility
and method construction. They do not yet replace the native event-contract
dependency inside the complete program runner. State invariants, opening
validation, report and carry composition, and preservation of stopped program
prefixes must retain their separate checks. General company rules and
independent implementation acceptance remain open.

\begingroup\raggedright
The ordered rules, full comparisons, identity repair and control results are
retained under
\path{data/accounting-master/ordered-event-eligibility-20261026/}.
\par
\endgroup
